% from ega0-22-fr.tex
\oldpage[0\textsubscript{IV-1}]{200}

dans $V''$ sont \emph{linéairement indépendantes sur $K'$}. Or, dans
$\Omega^1_{A''}\otimes_{A''}K'$, les $d_{A''}h_i\otimes1_{K'}$ sont les images
canoniques (20.5.11.2) des $h_i$ par l'injection
$V''\longrightarrow\Omega^1_{A''}\otimes_{A''}K'$ (cf. (22.4.6.1)), et cela
achève de prouver (22.4.7).
\end{proof}

On a en outre prouvé~:

\begin{corollary}[22.4.7.4]
\label{0.22.4.7.4-fr}
Pour que $\operatorname{rg}_K(V)=\operatorname{rg}_{K'}(V')$, il faut et il suffit que les éléments
$d_Af_i\otimes1_K$ $(1\leqslant i\leqslant r)$ soient linéairement
indépendants dans le $K$-espace vectoriel $\Omega^1_A\otimes_AK$.
\end{corollary}

\begin{remark}[22.4.8]
\label{0.22.4.8-fr}
On peut définir les homomorphismes (22.4.5.2) et (22.4.6.2) dans des conditions
un peu différentes. Soient $A$ un anneau, $\mathfrak m$ un idéal \emph{de carré
nul} dans $A$, $K=A/\mathfrak m$, $p$ un nombre premier tel que
$p\cdot1\in\mathfrak m$ (autrement dit $K$ est de caractéristique $p$, mais non
nécessairement $A$). Soit d'autre part $B$ une $A$-algèbre qui est un $A$-module
\emph{fidèlement plat} (de sorte que $A\subset B$), et supposons que l'on ait
\begin{equation}
\label{0.22.4.8.1-fr}
  y^p\in A+pB\subset A+\mathfrak mB
  \tag{22.4.8.1}
\end{equation}
pour tout $y\in B$. Si $y^p=x+pz$ avec $x\in A$, $z\in B$, la \emph{classe}
mod.~$(A\cap pB)$ de $x$ est bien déterminée, et comme $A\cap pB=pA$ puisque
$B$ est un $A$-module fidèlement plat, on a ainsi défini une application
canonique $j:B\longrightarrow A/pA$, telle que, pour $x\in A$, $j(x)$ soit la
classe de $x^p$ mod.~$pA$. D'ailleurs, si $y'\in B$ et
$y'^p=x'+pz'$ avec $x'\in A$, $z'\in B$, il est immédiat que les éléments
$(y+y')^p-x-x'$ et $(yy')^p-xx'$ appartiennent à $pB$ en vertu du fait que les
coefficients binomiaux $\binom pi$ sont des multiples de $p$ pour
$1\leqslant i\leqslant p-1$. L'application $j$ est donc un
\emph{homomorphisme d'anneaux}. Par composition, on déduit de $j$ un
homomorphisme $j':B\xrightarrow{j}A/pA\longrightarrow A/\mathfrak m=K$ et par
suite un homomorphisme $j_0:B_{(K)}\longrightarrow K$ tel que $j'$ soit le
composé $B\longrightarrow B_{(K)}=B\otimes_AK\xrightarrow{j_0}K$~; en outre,
si $i_0:K\longrightarrow B_{(K)}$ est l'application canonique, $i_0$ et $j_0$
vérifient les conditions de (21.3.1) pour les anneaux \emph{de caractéristique
$p$}, $K$ et $B_{(K)}$. Cela étant, pour définir un $K$-homomorphisme
$\pi'_{B_{(K)}/K}:\Omega^1_{B_{(K)}/K}\otimes_{B_{(K)}}K_{[j_0]}
\longrightarrow\Omega^1_A\otimes_AK$, il suffit (20.4.8) de définir une
$K$-dérivation
\begin{equation}
\label{0.22.4.8.2-fr}
  D_0:B_{(K)}\longrightarrow\Omega^1_A\otimes_AK
  \tag{22.4.8.2}
\end{equation}
où le second membre est considéré comme un $B_{(K)}$-module au moyen de $j_0$.
Pour cela, il suffit de définir une $A$-dérivation
\begin{equation}
\label{0.22.4.8.3-fr}
  D:B\longrightarrow\Omega^1_A\otimes_AK
  \tag{22.4.8.3}
\end{equation}
qui \emph{s'annule dans} $\mathfrak mB$ (le second membre étant considéré comme
$B$-module au moyen de $j'$). Or, si l'on a la relation $y^p=x+pz$ avec
$y\in B$, $x\in A$, $z\in B$, l'élément $d_Ax\otimes1_K$ de
$\Omega^1_A\otimes_AK$ ne dépend que de la classe $j(y)$ mod.~$pA$ de $x$
(donc seulement de $y$), car pour tout $t\in A$ on a
$d_A(pt)=p\,d_A(t)$, donc $d_A(pt)\otimes1=0$ dans
$\Omega^1_A\otimes_AK=\Omega^1_A/\mathfrak m\Omega^1_A$, puisque
$p\cdot1\in\mathfrak m$. On peut donc prendre
\[
  D(y)=d_Ax\otimes1_K.
\]
Le fait que $D$ est une dérivation résulte aisément de ce que $j$ est un
homomorphisme d'anneaux~; en outre, pour $t\in A$, on a
$d_A(t^px)=t^p d_A(x)+pt^{p-1}x d_A(t)$ et l'on en tire que
$D(ty)=tD(y)$, donc $D$ est une $A$-dérivation~; enfin, si l'on a de plus
$t\in\mathfrak m$, on a

